\documentclass[11pt,fleqn]{article}
\usepackage[utf8]{inputenc}
\usepackage{amsmath,amssymb,amsthm}
\usepackage{geometry}
\usepackage{booktabs}
\usepackage{graphicx}
\usepackage{microtype}
\usepackage{hyperref}
\usepackage{cite}

\geometry{letterpaper, margin=1in}

\hypersetup{
    colorlinks=true,
    linkcolor=blue,
    citecolor=blue,
    urlcolor=blue
}

\newtheorem{theorem}{Theorem}
\newtheorem{proposition}{Proposition}
\newtheorem{remark}{Remark}

\title{\textbf{Spurious Finite-Time Modal Blowup in Low-Order Galerkin Approximations of Fast--Slow Stably Stratified Boundary Layers:}\\
A PhD Research Project}
\author{\textbf{PhD Candidate:} \textit{[Student Name]} \\ \textbf{Supervisor:} David England (UAH.Edu BLM)}
\date{\today}

\begin{document}
\maketitle

\section{Project overview}

The stably stratified atmospheric boundary layer (SBL) exhibits fast--slow dynamics
between turbulent kinetic energy (TKE) and mean shear, with strong timescale
separation and sensitivity to stratification. Low-order Galerkin closures and
single-column models can display spurious finite-time blowup or collapse of
modal amplitudes, obscuring the physical interpretation of turbulence
intermittency and regime transitions.

This PhD project will develop a rigorous fast--slow dynamical-systems framework
for a prototypical SBL model, characterize its critical manifold and fold
structure, and relate these features to numerical pathologies in low-order
Galerkin approximations and operational boundary-layer schemes.

\section{Core model and fast--slow structure}

We consider a singularly perturbed system for turbulent kinetic energy $E$ and
mean shear $S$, with timescale separation $0 < \epsilon \ll 1$:
\begin{align}
\epsilon \frac{dE}{dt} &= f(E,S) = l \sqrt{E + \delta} \left( S^2 - \phi N^2 - c_b N^2 \frac{E}{E + \alpha} \right) - \frac{E^{3/2}}{l}, \label{eq:fast_tke} \\
\frac{dS}{dt} &= g(E,S) = G - c_1 E S - c_2 S, \label{eq:slow_shear}
\end{align}
on the physical domain $\mathcal{D} = \{(E,S) \in \mathbb{R}^2 : E \ge 0\}$.
Here $l$ is the mixing length, $N^2$ the (constant) Brunt--Väisälä frequency,
$G$ the geostrophic forcing, $\phi > 0$ a stability-correction coefficient,
$\delta, \alpha > 0$ regularization/saturation constants, $c_b = 0.50$, and
$c_1 = 1.80\ \text{s\,m}^{-2}$. The critical manifold is
\begin{equation}
\mathcal{C}_0 = \{(E,S) \in \mathcal{D} : f(E,S) = 0\}.
\end{equation}

\section{Research objectives}

\subsection{Objective 1: Well-posedness and regularity of the fast subsystem}

\begin{itemize}
    \item Establish dimensional consistency of $f(E,S)$ and $g(E,S)$, determining
    the required units of $\delta$, $\alpha$, $\phi$, and $c_2$ so that
    $f$ carries units of $\mathrm{m}^2\,\mathrm{s}^{-3}$ and $E/(E+\alpha)$ is
    dimensionless, with $[G] = \mathrm{s}^{-2}$.
    \item Analyze $C^1$-regularity of $f$ on $\mathcal{D}$, showing that
    $\delta>0$ restores differentiability at $E=0$ and yields a well-posed fast
    subsystem for geometric singular perturbation theory.
\end{itemize}

\subsection{Objective 2: Critical manifold geometry and fold structure}

\begin{itemize}
    \item Derive the fast eigenvalue
    \begin{equation}
        \lambda_f(E,S) \equiv \frac{\partial f}{\partial E}(E,S),
    \end{equation}
    in closed form, including the product rule on $\sqrt{E+\delta}$, the
    quotient/chain rule on $E/(E+\alpha)$, and the $E^{3/2}/l$ term.
    \item Apply the implicit function theorem to show that at points
    $(E^*,S^*) \in \mathcal{C}_0$ with $\lambda_f(E^*,S^*) \neq 0$, the set
    $\mathcal{C}_0$ is locally the graph of a $C^1$ function $E = E^*(S)$, and
    characterize the loss of regularity at $\lambda_f=0$.
    \item Classify branches of $\mathcal{C}_0$ into attracting ($\lambda_f<0$),
    repelling ($\lambda_f>0$), and fold ($\lambda_f=0$) segments, using the fast
    subsystem $\epsilon \dot E = f(E,S)$ with $S$ frozen to determine the
    timescale of attraction/repulsion.
\end{itemize}

\subsection{Objective 3: Nondegenerate folds and slow passage}

\begin{itemize}
    \item Identify candidate fold points $(E_0,S_0)$ satisfying
    $f(E_0,S_0)=0$ and $f_E(E_0,S_0)=0$, and compute $f_{EE}(E_0,S_0)$ and
    $f_S(E_0,S_0)$ explicitly.
    \item Formulate and verify nondegeneracy ($f_{EE}\neq 0$) and transversality
    ($f_S\neq 0$) conditions for generic saddle-node folds of the fast
    subsystem, interpreting each condition in terms of excluded degeneracies.
    \item For a nondegenerate fold with slow drift
    $\dot S \approx g(E_0,S_0) \equiv v_0 \neq 0$, approximate


\[
    f(E,S) \approx f_S(E_0,S_0)(S - S_0) + \tfrac12 f_{EE}(E_0,S_0)(E-E_0)^2,
    \]


    substitute $S-S_0 \approx v_0 t$ into $\epsilon \dot E = f(E,S)$, and
    derive the fractional scaling exponent relating the fold-crossing transient
    duration to $\epsilon$ (showing it is not $\mathcal{O}(\epsilon)$).
\end{itemize}

\subsection{Objective 4: Reduced slow flow, hysteresis, and Galerkin blowup}

\begin{itemize}
    \item On attracting branches, derive the reduced slow dynamics


\[
    \frac{dS}{dt} = g(E^*(S), S),
    \]


    where $E^*(S)$ solves $f(E,S)=0$ on the attracting portion of
    $\mathcal{C}_0$.
    \item Analyze multivalued (S-shaped) $E^*(S)$ in the presence of fold pairs,
    and show how slow variation of $S$ or forcing $G$ produces hysteresis in
    $E$, corresponding to abrupt collapse and re-establishment of turbulence in
    the SBL.
    \item Connect the fold geometry and slow passage phenomena to spurious
    finite-time modal blowup in low-order Galerkin approximations, identifying
    parameter regimes and truncation choices that amplify or suppress such
    pathologies.
\end{itemize}

\section{Expected contributions}

The thesis will provide:
\begin{itemize}
    \item A mathematically rigorous fast--slow analysis of a prototypical SBL
    closure with regularization.
    \item A classification of critical manifold branches, folds, and hysteresis
    mechanisms relevant to turbulence collapse and recovery.
    \item A diagnostic framework for detecting and mitigating spurious
    finite-time blowup in low-order Galerkin and single-column boundary-layer
    models.
\end{itemize}

\end{document}
